% Week 2 corrections included by Corrections.tex.

\section{Week 2 --- Brownian Motion, Stopping Times, and Quadratic Variation}

\section*{Part I --- Gaussian Structure and Brownian Transformations}

\subsection*{Exercise 2.1 --- Gaussian vectors and the covariance kernel}
\begin{enumerate}
  \item Independent increments give
  \[
  (W_r,W_s-W_r,W_t-W_s)
  \sim\mathcal N\!\left(
  0,\operatorname{diag}(r,s-r,t-s)
  \right).
  \]
  \item The level vector is a linear transformation of this Gaussian vector,
  hence is Gaussian, with covariance matrix
  \[
  \begin{pmatrix}
  r&r&r\\
  r&s&s\\
  r&s&t
  \end{pmatrix}.
  \]
  \item Since \(W_t=W_s+(W_t-W_s)\),
  \[
  W_t\mid\F_s\sim\mathcal N(W_s,t-s),
  \qquad
  W_t\mid(W_s=x)\sim\mathcal N(x,t-s).
  \]
  \item For any dates and coefficients,
  \[
  \sum_{i,j}a_ia_jK(t_i,t_j)
  =\operatorname{Var}\!\left(\sum_i a_iW_{t_i}\right)\geq0.
  \]
  \item Every finite vector of Brownian values is centered Gaussian. A
  Gaussian law is determined by its mean and covariance matrix, whose entries
  are generated by \(K\).
\end{enumerate}

\subsection*{Exercise 2.2 --- Symmetry, scaling, renewal, and time reversal}
\begin{enumerate}
  \item Both processes start from zero and are continuous. Their increments
  over disjoint intervals are independent, and
  \[
  X_t-X_u\sim\mathcal N(0,t-u),
  \qquad
  Y_t-Y_u=\frac{W_{ct}-W_{cu}}{\sqrt c}
  \sim\mathcal N(0,t-u).
  \]
  Hence both are standard Brownian motions.
  \item \(Z_0=0\), and its increments are Brownian increments after time
  \(s\). They are independent of one another and of \(\F_s\), Gaussian with
  the correct variances, and continuous.
  \item For \(0\leq u<t\leq s\),
  \[
  R_t-R_u=W_{s-t}-W_{s-u}\sim\mathcal N(0,t-u).
  \]
  Disjoint reversed intervals yield independent increments. Together with
  continuity, this proves equality of the finite-dimensional laws.
  \item Applying scaling to the whole path,
  \[
  (W_{cu})_{0\leq u\leq t}
  \overset{d}=(\sqrt c W_u)_{0\leq u\leq t}.
  \]
  Taking maxima gives the required identity.
  \item Almost surely,
  \[
  \begin{aligned}
  \limsup_{t\to\infty}\frac{W_t}{\sqrt{2t\log\log t}}&=1,
  &\liminf_{t\to\infty}\frac{W_t}{\sqrt{2t\log\log t}}&=-1,\\
  \limsup_{t\downarrow0}\frac{W_t}{\sqrt{2t\log\log(1/t)}}&=1,
  &\liminf_{t\downarrow0}\frac{W_t}{\sqrt{2t\log\log(1/t)}}&=-1.
  \end{aligned}
  \]
\end{enumerate}

\section*{Part II --- Markov Property, Martingales, and Stopping}

\subsection*{Exercise 2.3 --- Markov semigroup and martingale correction}
\begin{enumerate}
  \item Write \(W_{t+h}=W_t+G_h\), where
  \(G_h\sim\mathcal N(0,h)\) is independent of \(\F_t\). Conditioning gives
  \[
  \E[f(W_{t+h})\mid\F_t]=(P_hf)(W_t).
  \]
  \item Applying the Markov identity twice,
  \[
  \E[f(W_{t+h+k})\mid\F_t]
  =P_h(P_kf)(W_t)=P_{h+k}f(W_t),
  \]
  so \(P_{h+k}=P_hP_k\).
  \item Independent centered increments yield
  \[
  \E[W_t\mid\F_s]=W_s,
  \qquad
  \E[W_t^2\mid\F_s]=W_s^2+(t-s).
  \]
  \item The first identity proves that \(W\) is a martingale. The second
  gives
  \[
  \E[W_t^2-t\mid\F_s]=W_s^2-s,
  \]
  so \(W_t^2-t\) is also a martingale.
  \item For every continuous square-integrable martingale \(X\), there is a
  unique increasing predictable process \(\langle X\rangle\), starting from
  zero, such that \(X_t^2-\langle X\rangle_t\) is a martingale. For Brownian
  motion, \(\langle W\rangle_t=[W]_t=t\).
\end{enumerate}

\subsection*{Exercise 2.4 --- Stopping times and information at stopping}
\begin{enumerate}
  \item The deterministic time \(T\), the first hit \(\tau_a\), and the
  first exit \(\tau_{\ell,u}\) are stopping times. Continuity makes their
  occurrence detectable from the observed path. The last zero \(\rho_T\)
  and the time of the maximum \(\gamma_T\) require future information and
  are not stopping times.
  \item By continuity, the path hits \(a\) by time \(t\) exactly when its
  maximum over \([0,t]\) is at least \(a\).
  \item
  \[
  \F_\tau=\{A\in\F:A\cap\{\tau\leq t\}\in\F_t
  \text{ for every }t\geq0\}.
  \]
  For continuous adapted \(W\), the stopped variable \(W_\tau\) is obtained
  as a limit of values at simple stopping times, so it is
  \(\F_\tau\)-measurable.
  \item Strong Markov says that
  \[
  (W_{\tau+v}-W_\tau)_{v\geq0}
  \]
  is Brownian motion independent of \(\F_\tau\). At \(\tau_a\), the fresh
  process is \((W_{\tau_a+v}-a)_{v\geq0}\).
  \item The reflection starts at the random time \(\tau_a\). Deterministic
  Markov does not justify renewal there; strong Markov does, and symmetry then
  permits the sign flip of the fresh continuation.
\end{enumerate}

\section*{Part III --- Reflection and Hitting Times}

\subsection*{Exercise 2.5 --- Reflection principle and first passage}
\begin{enumerate}
  \item On \(\{\tau_a\leq t\}\), define
  \[
  W_u^*=\begin{cases}
  W_u,&u\leq\tau_a,\\
  2a-W_u,&u>\tau_a.
  \end{cases}
  \]
  This pairs \(\{M_t\geq a,W_t<a\}\) with \(\{W_t>a\}\).
  \item Hence, for \(a>0\),
  \[
  \Pb(M_t\geq a)=\Pb(\tau_a\leq t)
  =2\left(1-\Phi\!\left(\frac{a}{\sqrt t}\right)\right).
  \]
  Thus \(\Pb(M_t\leq a)=2\Phi(a/\sqrt t)-1\).
  \item
  \[
  f_{\tau_a}(t)=\frac{a}{\sqrt{2\pi t^3}}
  \exp\!\left(-\frac{a^2}{2t}\right),
  \qquad t>0.
  \]
  \item The tail of \(M_t\) equals that of \(|W_t|\), so
  \[
  M_t\overset{d}=|W_t|,
  \qquad
  \E[M_t]=\E[|W_t|]=\sqrt{\frac{2t}{\pi}}.
  \]
  \item By scaling,
  \[
  \Pb(\tau_a\leq t)
  =2\left(1-\Phi\!\left(\frac{a}{\sqrt t}\right)\right)
  =\Pb(a^2\tau_1\leq t).
  \]
  \item Letting \(t\to\infty\) in the cdf gives
  \(\Pb(\tau_a<\infty)=1\). Moreover,
  \[
  \Pb(\tau_a>t)
  =2\Phi(a/\sqrt t)-1
  \sim a\sqrt{\frac{2}{\pi t}},
  \]
  whose integral over large \(t\) diverges. Hence \(\E[\tau_a]=\infty\).
\end{enumerate}

\section*{Part IV --- Quadratic Variation}

\subsection*{Exercise 2.6 --- Realized quadratic variation}
\begin{enumerate}
  \item The increments are independent \(\mathcal N(0,T/n)\), so
  \[
  \E[Q_n]=T,
  \qquad
  \operatorname{Var}(Q_n)=n\,2(T/n)^2=\frac{2T^2}{n}.
  \]
  \item Therefore \(\E[(Q_n-T)^2]=2T^2/n\to0\).
  \item By the mean value theorem,
  \[
  \sum_i(\Delta_iA)^2
  \leq\|A'\|_\infty^2\,T\,|\pi_n|\longrightarrow0.
  \]
  \item Cauchy--Schwarz gives
  \[
  \left|\sum_i\Delta_iA\,\Delta_iW\right|
  \leq
  \left(\sum_i(\Delta_iA)^2\right)^{1/2}Q_n^{1/2}.
  \]
  The first factor tends to zero and \(Q_n\to T\), so the product tends to
  zero in probability.
  \item Expanding squared increments and using the preceding limits yields
  \([A+\sigma W]_T=\sigma^2T\).
  \item A consistent estimator is
  \[
  \widehat\sigma_n^2
  =\frac1T\sum_{i=0}^{n-1}(X_{t_{i+1}}-X_{t_i})^2
  \xrightarrow{\Pb}\sigma^2.
  \]
\end{enumerate}

\section*{Part V --- Price Models and Trading Gains}

\subsection*{Exercise 2.7 --- Symbolic analysis of geometric Brownian motion}
\begin{enumerate}
  \item
  \[
  \log(S_t/S_0)
  \sim\mathcal N\!\left((\mu-\tfrac12\sigma^2)t,\sigma^2t\right),
  \]
  and, conditionally on \(\F_s\),
  \[
  \log(S_t/S_s)
  \sim\mathcal N\!\left((\mu-\tfrac12\sigma^2)(t-s),
  \sigma^2(t-s)\right).
  \]
  \item Lognormal moments give
  \[
  \E[S_t]=S_0e^{\mu t},
  \qquad
  \operatorname{Var}(S_t)
  =S_0^2e^{2\mu t}(e^{\sigma^2t}-1),
  \]
  and
  \[
  \operatorname{Median}(S_t)
  =S_0e^{(\mu-\sigma^2/2)t}.
  \]
  \item
  \[
  \Pb(S_t>K)
  =1-\Phi\!\left(
  \frac{\log(K/S_0)-(\mu-\sigma^2/2)t}{\sigma\sqrt t}
  \right).
  \]
  \item Since \(\log S_t\) is a smooth drift plus \(\sigma W_t\),
  \([\log S]_t=\sigma^2t\). The drift has zero quadratic variation.
  \item If \(R_i=\log(S_{t_{i+1}}/S_{t_i})\), then
  \[
  \frac1T\sum_iR_i^2\xrightarrow{\Pb}\sigma^2
  \]
  as the mesh tends to zero.
\end{enumerate}

\subsection*{Exercise 2.8 --- A two-step adapted trading gain}
\begin{enumerate}
  \item At time \(t_1\), \(\Delta_1\) may use the path already observed but
  not the future increment \(W_T-W_{t_1}\).
  \item
  \[
  \E[\Delta_1(W_T-W_{t_1})]
  =\E[\Delta_1\E[W_T-W_{t_1}\mid\F_{t_1}]]=0.
  \]
  Also \(\E[W_{t_1}]=0\), hence \(\E[G_T]=0\).
  \item The cross-term is zero because
  \[
  \E[W_{t_1}\Delta_1(W_T-W_{t_1})]=0.
  \]
  Moreover,
  \[
  \E[\Delta_1^2(W_T-W_{t_1})^2]
  =(T-t_1)\E[\Delta_1^2].
  \]
  Therefore
  \[
  \E[G_T^2]=\Delta_0^2t_1+(T-t_1)\E[\Delta_1^2].
  \]
  \item The proposed position depends on the sign of a future increment and
  is not \(\F_{t_1}\)-measurable.
  \item Set \(\Delta_t=\Delta_0\) on \([0,t_1]\) and
  \(\Delta_t=\Delta_1\) on \((t_1,T]\). Then
  \[
  G_T=\int_0^T\Delta_t\dd W_t.
  \]
\end{enumerate}
